\chapter*{À propos de l'auteur}
\label{auteur}
\addcontentsline{toc}{chapter}{À propos de l'auteur}
\markboth{À propos de l'auteur}{À propos de l'auteur}

\textbf{OUEDRAOGO Ahmed Dieudonné} est médecin. C'est au cours de sa thèse de doctorat en médecine qu'il découvre Epi Info, en cherchant à analyser lui-même les données de son enquête. Ce premier contact avec l'analyse de données oriente durablement son parcours.

Data scientist, il s'intéresse à toutes les étapes du traitement des données de santé : conception des outils de collecte, nettoyage, analyse statistique, visualisation et automatisation des analyses. Il aime rendre ces méthodes accessibles au plus grand nombre, à commencer par des outils gratuits et simples d'utilisation comme Epi Info.

Consultant en e-santé, il accompagne la mise en place de solutions numériques pour la collecte et la gestion des données sanitaires, convaincu que des données de qualité sont la base de toute bonne décision en santé publique.

Il poursuit actuellement une formation d'épidémiologiste de terrain, centrée sur la surveillance épidémiologique, l'investigation des épidémies et la réponse aux urgences de santé publique. Sur le terrain, il mesure chaque jour combien un questionnaire bien conçu et une analyse rigoureuse comptent.

Ce livre est né de son expérience d'étudiant et de praticien. Il l'a écrit pour que d'autres étudiants, et tous les professionnels de santé qui débutent dans l'analyse de données, disposent du guide pratique en français qui lui avait manqué.

\vfill
\begin{center}
	\small
	© 2026 OUEDRAOGO Ahmed Dieudonné\\[0.5ex]
	Cet ouvrage est mis à disposition selon les termes de la licence\\
	\href{https://creativecommons.org/licenses/by-nc-sa/4.0/deed.fr}{Creative Commons Attribution -- Pas d'utilisation commerciale -- Partage dans les mêmes conditions 4.0 International (CC BY-NC-SA 4.0)}.\\[0.5ex]
	DOI : \href{https://doi.org/10.5281/zenodo.24182020}{10.5281/zenodo.24182020}
\end{center}
